\section{Details for the corner expansion}\label{sec:asym-details}

This appendix proves the estimates used in Sections~\ref{sec:corner}
and \ref{sec:asym}.  All constants in an $O$ term are independent
of the degree and of the row indices in their displayed ranges.

\subsection{Elementary series and quadrature facts}
\label{sec:asym-details-elementary}

We record the elementary cotangent expansion in the normalization used
here:
\begin{equation}\label{eq:asym-details-cot-fraction}
 \cot(\pi x)=\frac1\pi\left\{\frac1x+
       \sum_{k=1}^{\infty}\left(\frac1{x-k}+\frac1{x+k}\right)\right\}.
\end{equation}
The opposite poles are paired in this formula.  One direct proof is to
expand $\cos(xs)$ in its cosine series on $0\leq s\leq\pi$.
Integration gives its constant coefficient
$2\sin(\pi x)/(\pi x)$ and its $k$th coefficient
$2(-1)^kx\sin(\pi x)/(\pi(x^2-k^2))$.  These coefficients are
absolutely summable.  Evaluation of the uniformly convergent series
at $s=\pi$ gives \eqref{eq:asym-details-cot-fraction} for
nonintegral $x$.  The paired series and all its differentiated series
converge uniformly on compact sets away from the poles.  Comparing the
linear term at zero with the sine and cosine expansions gives
\begin{equation}\label{eq:asym-details-square-sums}
 \sum_{k\geq1}k^{-2}=\frac{\pi^2}{6},\qquad
 \sum_{\substack{k\geq1\\k\ \mathrm{odd}}}k^{-2}=\frac{\pi^2}{8}.
\end{equation}

We will repeatedly use the cellwise trapezoid identity
\begin{equation}\label{eq:asym-details-trapezoid-kernel}
 \frac{h}{2}\{g(u)+g(u+h)\}-\int_u^{u+h}g(t)\,dt
   =\frac12\int_u^{u+h}(t-u)(u+h-t)g''(t)\,dt.
\end{equation}
It follows by two integrations by parts.  For bounded $g''$ its
absolute value is at most $h^3\|g''\|_\infty/12$; it is also at
most $(h^2/8)\int_u^{u+h}|g''|$.  The corresponding midpoint error
is at most $h^3\|g''\|_\infty/24$, by integrating Taylor's formula
about the midpoint.

For $H_k=\sum_{j=1}^k j^{-1}$, applying
\eqref{eq:asym-details-trapezoid-kernel} to $1/t$ on the unit cells
from $k$ to infinity gives
\begin{equation}\label{eq:asym-details-harmonic-estimate}
 H_k=\log k+\gamma+\frac1{2k}+r_k,
 \qquad -\frac1{8k^2}\leq r_k\leq0.
\end{equation}
Indeed the sum of the nonnegative trapezoid errors in that tail is
$\log k+\gamma-H_k+1/(2k)$, and is bounded by
$\frac18\int_k^\infty 2t^{-3}\,dt=1/(8k^2)$.

The series \eqref{eq:corner-psi-definition} converges uniformly on
each compact subset of $(0,\infty)$: its $k$th summand is bounded
there by $C/k^2$.  Its derivative series have the same local uniform
convergence, and therefore
\begin{align}
 \psi(z+1)-\psi(z)&=\frac1z,
 &\psi'(z)&=\sum_{k=0}^{\infty}(z+k)^{-2},
 &\psi''(z)&=-2\sum_{k=0}^{\infty}(z+k)^{-3}.
 \label{eq:asym-details-psi-series}
\end{align}
These differentiations can also be justified by integrating the finite
derivative sums over a compact interval and passing to the uniform
limit.  In particular $\psi''<0$ is strictly increasing, and
\begin{equation}\label{eq:asym-details-psi-second-bound}
 |\psi''(z)|\leq z^{-2}+2z^{-3},\qquad z>0,
\end{equation}
by integral comparison of its positive series.

The same series gives
$\psi(z)=\lim_{K\to\infty}\{\log K-\sum_{k=0}^{K-1}(z+k)^{-1}\}$.
Apply \eqref{eq:asym-details-trapezoid-kernel} to
$(z+t)^{-1}$ on $[0,K]$ and then let $K\to\infty$; the bound
by the integral of its second derivative gives, uniformly for $z\geq1$,
\begin{equation}\label{eq:asym-details-psi-asymptotic}
 \psi(z)=\log z-\frac1{2z}+O(z^{-2}),\qquad
 \psi'(z)=\frac1z+O(z^{-2}).
\end{equation}
The second assertion also follows directly by integral comparison.
At the half integer the finite harmonic sums in the definition give
\begin{equation}\label{eq:asym-details-half-integer}
 \psi(\tfrac12)=-\gamma-2\log2,
 \qquad \psi'(\tfrac12)=4\sum_{k\ \mathrm{odd}}k^{-2}
                           =\frac{\pi^2}{2}.
\end{equation}
For example the finite sum for the first value is
$-\gamma+2H_K-2H_{2K}$, which has the indicated limit.

\subsection{The reference square and its defect}
\label{sec:asym-details-reference}

\begin{proof}[Proof of Lemma~\ref{lem:corner-reference-square}]
Simultaneous reflection $j\mapsto n-j$, $k\mapsto n-k$
preserves the denominator and the weights of the equal-mesh sum, and
replaces $1+\cos\theta_k$ by $1-\cos\theta_k$.  Averaging
therefore replaces its numerator by $1$.  There is no diagonal,
since $j+k$ is odd.  Put $h=\pi/(2n)$, $c_r=\csc(rh)$.
For $j>k$, set $s=j-k$, $r=j+k$.  Both are positive odd integers,
and
\[
 |\cos\theta_j-\cos\theta_k|=2\sin(rh)\sin(sh),
 \qquad s\leq r\leq2n-s.
\]
The two orders of $j,k$ cancel this factor $2$.  For fixed $s<n$,
the extreme values $r=s,2n-s$ have half weight, and the other
values have weight $1$.  Using $c_{2n-r}=c_r$, their sum is
\[
 c_s+2\sum_{\substack{s<r<n\\r\ \mathrm{odd}}}c_r
      +\boldsymbol1_{\{n\ \mathrm{odd}\}}c_n.
\]
When $s=n$ occurs, its sole pair $(j,k)=(n,0)$ has weight $1/4$.
Multiplying by $c_s$ and summing over $s$ is exactly the expansion
of $(\sum_{r\leq n,\ r\ \mathrm{odd}}d_rc_r)^2$.  Restoring
the factor $4/n^2$ proves the identity, including the odd endpoint.
\end{proof}

\begin{proof}[Proof of Lemma~\ref{lem:corner-one-dimensional}]
Set $y_r=(2r-1)\pi/(4n)$, $1\leq r\leq n$.  Pairing $r$ with
$n+1-r$ and using $\cot y+\tan y=2\csc(2y)$ gives
\begin{equation}\label{eq:asym-details-L-midpoints}
 L_n=\frac1n\sum_{r=1}^n\cot y_r.
\end{equation}
For odd $n$, the central midpoint is $\pi/4$ and contributes $1$,
which is precisely the half-weight endpoint in
\eqref{eq:corner-L-definition}.

Put $v_0(y)=\cot y-1/y$, $v_0(0)=0$, on $[0,\pi/2]$.
Then $\int_0^{\pi/2}v_0=\log(2/\pi)$.  The partial fraction
\eqref{eq:asym-details-cot-fraction}, rescaled to the variable $y$,
gives
\[
 |v_0''(y)|=4y\sum_{k\geq1}
       \frac{3\pi^2k^2+y^2}{(\pi^2k^2-y^2)^3}
 \leq\frac{416}{27\pi^3}\sum_{k\geq1}k^{-4}<1.
\]
For the last inequality use $\pi>3$ and
$\sum k^{-4}\leq1+\int_1^\infty t^{-4}dt=4/3$.
The composite midpoint estimate on these $n$ cells consequently gives
\[
 \left|\frac1n\sum_{r=1}^n v_0(y_r)
       -\frac2\pi\log(2/\pi)\right|
       \leq\frac{\pi^2}{96n^2}<\frac1{8n^2}.
\]
The singular part of \eqref{eq:asym-details-L-midpoints} is exactly
$a(2H_{2n}-H_n)$.  Formula
\eqref{eq:asym-details-harmonic-estimate} gives
\[
 |2H_{2n}-H_n-(\log n+\gamma+2\log2)|
       \leq\frac3{16n^2}.
\]
Its contribution after multiplication by $a=2/\pi$ is less than
$1/(8n^2)$.  Adding the two errors proves
\eqref{eq:corner-L-estimate}.
\end{proof}

To control the mesh deformation, we need a bound which is uniform up
to both ends of every diagonal.

\begin{lemma}[A uniform deformation bound]
\label{lem:asym-details-deformation}
For $0\leq j,k\leq n$, $j+k$ odd and $\ell=|j-k|$,
\begin{equation}\label{eq:asym-details-deformation-bound}
 |f(\theta_j,\phi_k)-f(\theta_j,\theta_k)|
       \leq\frac{\pi^2n}{2\ell^2}.
\end{equation}
\end{lemma}

\begin{proof}
There is no diagonal.  The cases $k=0$ and $j=n$ have zero
deformation, so suppose $k>0$, $j<n$.  Put
$u=\theta_j$, $v_0=\theta_k$, $v_1=\phi_k$,
$\delta=v_0-v_1=\pi k/(nm)$, $V=(v_0+v_1)/2$,
$S_i=(u+v_i)/2$, and $Q_i=|u-v_i|/2$.  The order of $u,v_i$
is the same for $i=0,1$.  Exact subtraction gives
\[
 E:=|f(u,v_1)-f(u,v_0)|
 =\frac{\cos^2(u/2)\sin V\sin(\delta/2)}
        {\sin S_0\sin S_1\sin Q_0\sin Q_1}.
\]
Since $\sin S_i\geq\cos(u/2)\sin(v_i/2)$,
\[
 \frac{\cos^2(u/2)\sin V}{\sin S_0\sin S_1}
 \leq\cot(v_0/2)+\cot(v_1/2)
 \leq2\cot(v_1/2)\leq\pi\frac{m-k}{k}.
\]
The last inequality uses $\sin z\geq2z/\pi$ and
$\cos z\leq\pi/2-z$ for $0\leq z\leq\pi/2$.
Multiplication by $\sin(\delta/2)\leq\pi k/(2nm)$ bounds
the numerator after this comparison by $\pi^2(m-k)/(2nm)$.

If $k<j$, both $\sin Q_i\geq\ell/n$, giving the result.
If $k>j$, write $r=n-j$.  The same chord inequality gives
\[
 \sin Q_0\geq\ell/n,\qquad
 \sin Q_1\geq\frac{n(\ell-1)+r}{nm},\qquad m-k\leq r.
\]
Thus
\[
 E\leq\frac{\pi^2nr}{2\ell\{n(\ell-1)+r\}}
       \leq\frac{\pi^2n}{2\ell^2},
\]
where the last step uses $r\leq n$.  This also covers
$\ell=1$ and $k=n$.
\end{proof}

\begin{proof}[Proof of Lemma~\ref{lem:corner-defect-limit}]
Using the actual weights in \eqref{eq:corner-positive-sum} and
Lemma~\ref{lem:corner-reference-square}, set
\[
 \Delta_n=\frac4{nm}
 \sum_{\substack{0\leq j,k\leq n\\j+k\ \mathrm{odd}}}
 \varepsilon_j\eta_k
      \{f(\theta_j,\phi_k)-f(\theta_j,\theta_k)\}.
\]
The $k=n$ reference row is zero, so its weight may be changed from
$\eta_n=1$ to $\varepsilon_n=1/2$ without changing its value.
Consequently the exact normalization is
\begin{equation}\label{eq:asym-details-exact-defect}
 \Lambda_c(n)-L_n^2=\Delta_n+\frac{1-L_n^2}{m}.
\end{equation}
The elementary bound $\sin(r\pi/(2n))\geq r/n$ shows
$L_n\leq2\sum_{r\leq n,\ r\ \mathrm{odd}}1/r\leq2+\log n$;
the second term in \eqref{eq:asym-details-exact-defect} therefore
tends to zero.

Fix an odd $\ell\geq1$ and an orientation $\sigma=\pm1$, and
write $k=j+\sigma\ell$, $t=j/n$.  On every compact subinterval
of $0<t<1$, expansion of the sine factors in the cosine difference
gives, uniformly as $n\to\infty$,
\begin{equation}\label{eq:asym-details-diagonal-limit}
 \frac{f(\theta_j,\phi_k)-f(\theta_j,\theta_k)}{n}
 \longrightarrow h_{\sigma,\ell}(t)
 =\frac{\cot(\pi t/2)}\pi
       \left(\frac1{|\sigma\ell-t|}-\frac1\ell\right).
\end{equation}
Indeed $\phi_k-\theta_j=\pi(\sigma\ell-t)/m$ and
$\theta_k-\theta_j=\pi\sigma\ell/n$; the other sine factor
tends to $\sin(\pi t)$, and
$(1+\cos(\pi t))/\sin(\pi t)=\cot(\pi t/2)$.
Taking the limit in
\eqref{eq:asym-details-deformation-bound} gives
$|h_{\sigma,\ell}|\leq\pi^2/(2\ell^2)$ in the open interval.
The explicit expression has finite endpoint limits.

For completeness, this convergence suffices for the full sums, rather
than only for interior rows.  For each fixed $\ell$, restrict first
to $\delta\leq j/n\leq1-\delta$.  Uniform convergence in
\eqref{eq:asym-details-diagonal-limit} and the usual Riemann sum
limit apply there; all weights are one for sufficiently large $n$.
The at most $2\delta n+O(1)$ remaining rows have total normalized
absolute contribution at most $C(\delta+1/n)/\ell^2$, by
\eqref{eq:asym-details-deformation-bound}.  The same bound controls
the omitted endpoint intervals of the limiting integral.  Let
$n\to\infty$, and then $\delta\downarrow0$.  This proves the
limit for every fixed offset, with the factor $4n/m\to4$.
Finally the combined finite and limiting tails for $\ell>M$ are
bounded by $C\sum_{\ell>M}\ell^{-2}\to0$, uniformly in $n$.
This justifies both limit exchanges without a parity restriction or
an unestimated endpoint strip.

Pairing orientations in the limit gives
\[
 4(h_{+,\ell}+h_{-,\ell})
 =\frac8\pi\cot\frac{\pi t}{2}
       \frac{t^2}{\ell(\ell^2-t^2)}.
\]
Summing proves \eqref{eq:corner-kappa-limit}.  The terms with
$\ell\geq3$ converge uniformly with all endpoint values, and the
single $\ell=1$ term has its pole canceled by the cotangent zero
at $t=1$.  At zero the factor $t^2\cot(\pi t/2)$ tends to zero.
Thus the claimed continuous endpoint values also follow.
\end{proof}

\subsection{Uniform centering of the harmonic rows}
\label{sec:asym-details-centering}

\begin{lemma}[Uniform centering, including the end rows]
\label{lem:asym-details-centering}
Let $m\geq2$, $1\leq s\leq m-1$, $r=m-s$, and $q=s/(2m)$.
Then the row $D_{m,r}$ of \eqref{eq:corner-harmonic-row} satisfies
\begin{align}
 D_{m,r}={}&\psi(r+\tfrac12)+\psi(s)-\psi(q)-\psi(m)
                       -\frac{J_{m,r}}{2m}-\rho_{m,r},
 \label{eq:asym-details-centering-exact}\\
 J_{m,r}={}&r\psi'(r+\tfrac12)+s\psi'(s)-s\psi'(m),
 &0<\rho_{m,r}\leq\frac{257}{216m^2}<\frac5{4m^2}.
 \label{eq:asym-details-centering-error}
\end{align}
In particular the actual weighted remainder is
$O(\log m/m^2)$ after multiplication by $4/(\pi m)$ and
summation with the weights $w(r/m)$.
\end{lemma}

\begin{proof}
For $z>h>0$ Taylor's formula in integral form is
\[
 \psi(z-h)=\psi(z)-h\psi'(z)+\mathcal R(z,h),\qquad
 \mathcal R(z,h)=\int_0^h(h-t)\psi''(z-t)\,dt<0.
\]
By \eqref{eq:asym-details-psi-second-bound},
\begin{equation}\label{eq:asym-details-Taylor-bound}
 0<-\mathcal R(z,h)
       \leq\frac{h^2}{2}\{(z-h)^{-2}+2(z-h)^{-3}\}.
\end{equation}
The three shifted arguments in $D_{m,r}$ are exactly
\[
 r+q=r+\tfrac12-\frac r{2m},\qquad
 s-q=s-\frac s{2m},\qquad m-q=m-\frac s{2m}.
\]
Apply this Taylor identity at
$(z,h)=(r+1/2,r/(2m)),(s,q),(m,q)$, keeping $-\psi(q)$
unchanged.  This gives \eqref{eq:asym-details-centering-exact},
with
\[
 -\rho_{m,r}=\mathcal R(r+\tfrac12,r/(2m))
                         +\mathcal R(s,q)-\mathcal R(m,q).
\]
Since $\psi''$ is strictly increasing and the last two remainders
have the same shift, $\mathcal R(s,q)-\mathcal R(m,q)<0$.
Thus $\rho_{m,r}>0$, and
\[
 \rho_{m,r}\leq-\mathcal R(r+\tfrac12,r/(2m))-\mathcal R(s,q).
\]
The lower arguments are $r+q\geq r$ and
$s-q\geq3s/4$, respectively.  Formula
\eqref{eq:asym-details-Taylor-bound} therefore gives
\[
 -\mathcal R(r+\tfrac12,r/(2m))\leq\frac3{8m^2},\qquad
 -\mathcal R(s,q)\leq\frac{22}{27m^2}.
\]
Their sum is $257/(216m^2)$, proving the asserted uniform bound.

Finally $0<w(r/m)<am/r$ follows from $\tan z>z$ on
$(0,\pi/2)$.  Hence $\sum_{r<m}w(r/m)<amH_{m-1}$.
More precisely the weighted difference between the centered expression
and $\mathcal H_m$ is positive and at most
\[
 \frac4{\pi m}\frac{257}{216m^2}\,amH_{m-1}
       =\frac{257H_{m-1}}{27\pi^2m^2}.
\]
Since $H_{m-1}\leq1+\log(m-1)$, this proves the final assertion.
\end{proof}

\subsection{Aggregation of the harmonic triangle}
\label{sec:asym-details-harmonic}

We now prove \eqref{eq:asym-harmonic-expansion}.  Set
\[
 d_0=\gamma+2\log2,\qquad q_0=\frac{\pi^2}{4}-1,
 \qquad
 F(x)=\log x+\log(1-x)-\psi((1-x)/2).
\]
The auxiliary constants $d_0,q_0$ here are distinct from the final
coefficients $c_1,b,c$ in \eqref{eq:asym-coefficients}.

\paragraph{The summable endpoint corrections.}
For real $t\geq1$ define
\[
 e(t)=\psi(t+\tfrac12)-\log t,\qquad
 u(t)=t\psi'(t+\tfrac12)-1.
\]
The midpoint estimate for $y^{-2}$ on the unit intervals covering
$[t,\infty)$ gives
$\psi'(t+1/2)=1/t+O(t^{-3})$ uniformly for $t\geq1$:
the individual errors are bounded by a constant times $(t+k)^{-4}$,
whose sum is $O(t^{-3})$.  Consequently $u(t)=O(t^{-2})$.
Formula \eqref{eq:asym-details-psi-asymptotic} implies
$e(t)\to0$; integrating $e'(t)=O(t^{-3})$ from infinity gives
$e(t)=O(t^{-2})$.  Thus
\begin{equation}\label{eq:asym-details-endpoint-sums}
 S_e=\sum_{r=1}^{\infty}\frac{e(r)}r,
 \qquad S_u=\sum_{r=1}^{\infty}\frac{u(r)}r
\end{equation}
are absolutely convergent.

Only the second must be evaluated.  Finite recurrence summation gives
the exact identity
\begin{equation}\label{eq:asym-details-trigamma-finite-sum}
 \sum_{r=1}^R\psi'(r+\tfrac12)
  =(R+\tfrac12)\psi'(R+\tfrac12)+\psi(R+\tfrac12)
                         -\psi(\tfrac12)-\tfrac12\psi'(\tfrac12).
\end{equation}
To verify it, write
$\psi'(r+1/2)=\psi'(1/2)-\sum_{k=0}^{r-1}(k+1/2)^{-2}$,
sum this finite triangle, and use
$k=(k+1/2)-1/2$ in the coefficient of each summand.
Subtract $H_R$ in \eqref{eq:asym-details-trigamma-finite-sum}
and let $R\to\infty$.  The first term on the right tends to $1$,
and $\psi(R+1/2)-H_R\to-\gamma$.  In conjunction with
\eqref{eq:asym-details-half-integer}, this yields
\begin{equation}\label{eq:asym-details-Su-value}
 S_u=1+2\log2-\frac{\pi^2}{4},\qquad q_0+S_u=2\log2.
\end{equation}

\paragraph{An expansion that can be summed uniformly.}
Put $s=m-r$, $x=r/m$.  Substitute
\eqref{eq:asym-details-psi-asymptotic} into the centering lemma,
but retain $e(r)$ and $u(r)$ exactly.  The result is
\begin{align}
 D_{m,r}={}&\log m+F(x)+e(r)
       -\frac1{2m(1-x)}-\frac{x}{2m}-\frac{u(r)}{2m}
       +Q_{m,r},\label{eq:asym-details-row-expanded}\\
 |Q_{m,r}|\leq{}& C\left(s^{-2}+(ms)^{-1}+m^{-2}\right).
 \label{eq:asym-details-row-remainder}
\end{align}
Here the integer difference $\psi(s)-\psi(m)$ contributes
$\log(s/m)-1/(2s)+1/(2m)+O(s^{-2}+m^{-2})$.
The derivative bracket in \eqref{eq:asym-details-centering-exact}
is $1+x+u(r)+O(1/s+1/m)$; this cancels the scalar $1/(2m)$
and explains the term $-x/(2m)$ in
\eqref{eq:asym-details-row-expanded}.

For summation we have both endpoint bounds
\begin{equation}\label{eq:asym-details-weight-bounds}
 0<w(r/m)\leq\frac{am}{r},\qquad
 w(r/m)\leq C\frac{s}{m}\quad(r\geq m/2).
\end{equation}
The second follows from $w(r/m)=\tan(\pi s/(2m))$ on
$0\leq s/m\leq1/2$.  Split the sum at $r=m/2$.  On its first
half, $s\geq m/2$ and
$(1/m)w|Q|\leq C/(m^2r)$.  On the second half the $s^{-2}$
part is bounded by $C/(m^2s)$ after weighting and normalization;
the other two parts are no larger in total.  Hence
\begin{equation}\label{eq:asym-details-weighted-row-error}
 \frac1m\sum_{r=1}^{m-1}w(r/m)|Q_{m,r}|
          =O(\log m/m^2).
\end{equation}

Let $v(x)=w(x)-a/x$, with $v(0)=0$.  This is smooth on $[0,1]$
and $|v(x)|\leq Cx$.  The decay of $e,u$ and their series tails
in \eqref{eq:asym-details-endpoint-sums} give
\begin{align}
 \frac1m\sum_{r<m}w(r/m)e(r)&=aS_e+O(\log m/m^2),
 \label{eq:asym-details-e-weighted}\\
 \frac1m\sum_{r<m}w(r/m)u(r)&=aS_u+O(\log m/m^2).
 \label{eq:asym-details-u-weighted}
\end{align}
Indeed the tails of $\sum e(r)/r$, $\sum u(r)/r$ are
$O(m^{-2})$, and each $v$ contribution is bounded by
$Cm^{-2}\sum_{r<m}r^{-1}$.  In particular the $e$ term contributes
a constant but no $1/m$ term.

The functions $w(x)/(1-x)-a/x$ and $xw(x)$ are smooth through
both endpoints.  An ordinary trapezoid estimate and
\eqref{eq:asym-details-harmonic-estimate} therefore give
\begin{align}
 \frac1m\sum_{r<m}\frac{w(r/m)}{1-r/m}
       &=a\log m+a\gamma+I_1+O(1/m),
 \label{eq:asym-details-I1-sum}\\
 \frac1m\sum_{r<m}\frac rm w(r/m)&=I_2+O(1/m).
 \label{eq:asym-details-I2-sum}
\end{align}
If
\[
 M_m=\frac1m\sum_{r<m}w(r/m)\{\log m+F(r/m)\},
\]
the preceding formulas applied to
\eqref{eq:asym-details-row-expanded} now yield
\begin{equation}\label{eq:asym-details-H-to-M}
 \mathcal H_m=\frac4\pi\left\{
 M_m+aS_e-\frac{a\log m+a\gamma+I_1+I_2+aS_u}{2m}
                         \right\}+O(\log m/m^2).
\end{equation}

\paragraph{The logarithmic quadrature.}
Subtract the explicit singularity by setting
\begin{equation}\label{eq:asym-details-R-definition}
 R(x)=w(x)F(x)-a\frac{\log x+d_0}{x}.
\end{equation}
Near zero put $J(x)=\log(1-x)-\psi((1-x)/2)$.
It is analytic there with $J(0)=d_0$ and $J'(0)=q_0$, by
\eqref{eq:asym-details-half-integer}.  Hence
\[
 R(x)=a\frac{J(x)-d_0}{x}+v(x)\{\log x+J(x)\}.
\]
The first summand is analytic and $v(x)=O(x)$.  Near the other
endpoint put $y=1-x$.  The recurrence gives
$\psi(y/2)=-2/y+\psi(1+y/2)$, with its second term analytic,
whereas $w(1-y)=\tan(\pi y/2)$.  These actual analytic
factorizations, followed by compactness on an interior interval, prove
\begin{align}
 R(0)&=aq_0,&R(1)&=\pi-ad_0,
 \label{eq:asym-details-R-endpoints}\\
 |R''(x)|&\leq C\{x^{-1}+(1-x)^{-1}\},&&0<x<1,
 \label{eq:asym-details-R-second}\\
 |R(x)-R(0)|&\leq Cx(1+|\log x|),&&0<x\leq\tfrac12,
 \notag\\
 |R(1-y)-R(1)|&\leq Cy(1+|\log y|),&&0<y\leq\tfrac12.
 \label{eq:asym-details-R-endpoint-moduli}
\end{align}
In particular no differentiation of an unspecified $O$ term is needed
to obtain \eqref{eq:asym-details-R-second}.

Let $h=1/m$.  On each endpoint cell,
\eqref{eq:asym-details-R-endpoint-moduli} bounds the trapezoid error
by $O(h^2(1+|\log h|))$.  On the interior cells use the integral
bound in \eqref{eq:asym-details-trapezoid-kernel} and integrate
\eqref{eq:asym-details-R-second} over $[h,1-h]$.  This also costs
$O(h^2\log m)$.  Thus the actual endpoint-excluded formula is
\begin{equation}\label{eq:asym-details-R-quadrature}
 \frac1m\sum_{r<m}R(r/m)=\int_0^1R(x)\,dx
       -\frac{R(0)+R(1)}{2m}+O(\log m/m^2).
\end{equation}

Two harmonic sums finish the calculation:
\begin{align}
 H_{m-1}&=\log m+\gamma-\frac1{2m}+O(m^{-2}),
 \label{eq:asym-details-Hm1}\\
 \sum_{r<m}\frac{\log r}{r}
   &=\frac12\log^2m+C_K-\frac{\log m}{2m}
                                     +O(\log m/m^2).
 \label{eq:asym-details-log-harmonic}
\end{align}
The first follows from \eqref{eq:asym-details-harmonic-estimate}.
For the second apply the unit-cell trapezoid identity to
$g(t)=\log t/t$ on $[1,m]$.  The sum of its local errors converges,
since $g''(t)=(2\log t-3)/t^3$ is integrable at infinity, and its
tail is $O(\log m/m^2)$.  This proves the formula with a real
constant $C_K$.

Finally $v(0)=0$, $v(1)=-a$, and
$\int_0^1v(x)\,dx=a\log(2/\pi)$.  The ordinary trapezoid
formula combined with \eqref{eq:asym-details-Hm1} shows
\begin{equation}\label{eq:asym-details-w-sum}
 \frac1m\sum_{r<m}w(r/m)
       =a(\log m+\gamma)+a\log(2/\pi)+O(m^{-2}).
\end{equation}
Its $1/m$ terms cancel: $aH_{m-1}$ contributes $-a/(2m)$,
and the endpoint correction for $v$ contributes $+a/(2m)$.

Use \eqref{eq:asym-details-R-definition} to write $M_m$ exactly
as the sum of $\log m$ times the left side of
\eqref{eq:asym-details-w-sum}, the singular sum
\[
 a\sum_{r<m}\frac{\log r}{r}+a(d_0-\log m)H_{m-1},
\]
and $(1/m)\sum R(r/m)$.  Inserting
\eqref{eq:asym-details-R-quadrature}--\eqref{eq:asym-details-w-sum}
gives a real constant $C_M$ with
\begin{equation}\label{eq:asym-details-M-expansion}
 M_m=\frac a2\log^2m+ad\log m+C_M
       -\frac{\pi+aq_0}{2m}+O(\log m/m^2).
\end{equation}
Here $d=d_0+\log(2/\pi)$.  To see the first-order cancellation
explicitly, the $-a\log m/(2m)$ from
\eqref{eq:asym-details-log-harmonic} cancels the
$+a\log m/(2m)$ from $-a\log m H_{m-1}$.
The remaining coefficient is
$-(ad_0+R(0)+R(1))/2=-(\pi+aq_0)/2$.

Substitute \eqref{eq:asym-details-M-expansion} and
\eqref{eq:asym-details-Su-value} into
\eqref{eq:asym-details-H-to-M}.  Since $4/\pi=2a$ and
$q_0+S_u=2\log2$, the logarithmic first correction is
$-a^2\log m/m$, and its constant companion is
\[
 -a\{\pi+aq_0+a\gamma+I_1+I_2+aS_u\}
       =-2-a^2(\gamma+2\log2)-a(I_1+I_2)=\beta_H.
\]
All terms independent of $m$ can be collected in $C_H$ after
subtracting the constant term in $A(m)^2$.  This proves
\eqref{eq:asym-harmonic-expansion} with its stated uniform error.

\subsection{Quadrature for the smooth residual}
\label{sec:asym-details-smooth}

We finish the proof of Proposition~\ref{prop:asym-two-expansions}
by proving \eqref{eq:asym-smooth-expansion}, including its explicit
error and the nonvanishing first coefficient.

\paragraph{Uniform smoothness after removing both poles.}
Removing the poles $0$ and $1$ in
\eqref{eq:asym-details-cot-fraction} gives
\begin{equation}\label{eq:asym-details-B-series}
 B(x)=\frac1\pi\left\{\frac1{x+1}
       +\sum_{k\geq2}\left(\frac1{x-k}+\frac1{x+k}\right)\right\}.
\end{equation}
It is analytic on $[-1/4,5/4]$.  The remaining negative poles are
$-h$ and the positive ones $h+1$, $h\geq1$, both at distance
at least $3h/4$ from that interval.  Thus, for $j=1,2$,
\[
 |B^{(j)}(x)|\leq\frac{2j!}{\pi}(4/3)^{j+1}
                    \sum_{h\geq1}h^{-j-1}.
\]
Using $\sum h^{-2}\leq2$, $\sum h^{-3}\leq3/2$, and
$\pi>3$, we obtain the convenient bounds
\begin{equation}\label{eq:asym-details-B-bounds}
 |B'|\leq3,\qquad |B''|\leq5,\qquad |B|<5
                   \quad\hbox{on }[-1/4,5/4].
\end{equation}
The last bound follows from $B(0)=1/\pi$ and the first derivative
bound.  Also $B(1-x)=-B(x)$.  Consequently, if
\[
 I(t)=\int_0^t B(x)\,dx,
\]
then $I(0)=I(1)=0$, and direct integration gives
\begin{equation}\label{eq:asym-details-I-log}
 I(t)=\frac1\pi\log\frac{\sin(\pi t)}{\pi t(1-t)},
                      \qquad 0<t<1.
\end{equation}
The integral definition supplies its smooth endpoint values.

Put $W(t)=tw(t)$, $W(0)=a$.  We have $0\leq W\leq a<1$.
Substituting $x=t/2$ into
\eqref{eq:asym-details-cot-fraction} yields
\[
 W(t)=\frac2\pi-\frac{4t^2}{\pi}
                     \sum_{k\geq1}\frac1{4k^2-t^2}.
\]
Differentiating the uniformly convergent series on $[0,1]$, and using
$4k^2-t^2\geq3k^2$, gives
\begin{align*}
 |W'|&\leq\frac{32}{9\pi}\sum_{k\geq1}k^{-2}<3,\\
 |W''|&\leq\frac8\pi\left\{\frac49\sum_{k\geq1}k^{-2}
                       +\frac{16}{27}\sum_{k\geq1}k^{-4}\right\}<5.
\end{align*}
The elementary integral comparisons for these sums suffice for both
strict bounds.

\paragraph{The exact cell union in a row.}
Fix $t=r/m$ and $\alpha=(1-t)/2$.  The $r$ points
$(\ell+\alpha)/n$, $0\leq\ell<r$, are the midpoints of the
cells with endpoints $(\ell+\alpha\pm1/2)/n$.  Their union is
exactly
\begin{equation}\label{eq:asym-details-row-cell-union}
 [-t/(2n),\ t+t/(2n)].
\end{equation}
For $n\geq2$ this lies in $[-1/4,5/4]$, including the case
$r=n$.  The midpoint bound and \eqref{eq:asym-details-B-bounds}
give an error at most $5r/(24n^2)\leq5t/(16n)$ when comparing
the sum of $B$ values with $n$ times its integral on this union.
Put $\delta=t/(2n)$.  Each of the two extensions of $[0,t]$
satisfies
\[
 \left|n\int_{-\delta}^{0}B(x)\,dx-\frac t2B(0)\right|
       \leq\frac{3t^2}{8n},\qquad
 \left|n\int_t^{t+\delta}B(x)\,dx-\frac t2B(t)\right|
       \leq\frac{3t^2}{8n}.
\]
It follows that, for every row,
\begin{equation}\label{eq:asym-details-smooth-row}
 \sum_{\ell=0}^{r-1}B((\ell+\alpha)/n)
     =nI(t)+\frac t2\{B(t)+B(0)\}+e_{n,r},
 \qquad |e_{n,r}|\leq\frac{17t}{16n}.
\end{equation}

\paragraph{The outer quadrature and its normalization.}
Define smooth functions on $[0,1]$ by
\[
 \begin{aligned}
 J_B(t)&=\int_0^1B(st)\,ds,\\
 F_B(t)&=W(t)J_B(t)=w(t)I(t),\\
 G_B(t)&=W(t)\{B(t)+1/\pi\}.
 \end{aligned}
\]
The integral expression removes the apparent pole in $w(t)I(t)$.
From the preceding bounds,
$|J_B|\leq5$, $|J_B'|\leq3/2$, $|J_B''|\leq5/3$, so
\begin{equation}\label{eq:asym-details-FG-bounds}
 \|F_B''\|_\infty<36,\qquad\|G_B''\|_\infty<50.
\end{equation}
Their endpoint values are
\[
 F_B(0)=2/\pi^2,\quad F_B(1)=0,\qquad
 G_B(0)=4/\pi^2,\quad G_B(1)=0.
\]
Multiply \eqref{eq:asym-details-smooth-row} by its exact weight
$4w(t)/(nm)$ and sum.  Since $tw(t)\leq1$, this gives
\begin{equation}\label{eq:asym-details-smooth-outer}
 \frac{4V_n}{nm}=\frac4m\sum_{r=1}^{m-1}F_B(r/m)
          +\frac2{nm}\sum_{r=1}^{m-1}G_B(r/m)+E_n,
 \qquad |E_n|\leq\frac{17}{4nm}.
\end{equation}
The composite trapezoid formula and
\eqref{eq:asym-details-FG-bounds} give
\begin{align*}
 \frac4m\sum_{r=1}^{m-1}F_B(r/m)
   &=4\int_0^1F_B-\frac4{\pi^2m}+R_F,
       &|R_F|&\leq\frac{12}{m^2},\\
 \frac2{nm}\sum_{r=1}^{m-1}G_B(r/m)
   &=\frac2n\int_0^1G_B-\frac4{\pi^2nm}+R_G,
       &|R_G|&\leq\frac{25}{3nm^2}.
\end{align*}
The two endpoint corrections combine exactly:
\[
 -\frac4{\pi^2m}-\frac4{\pi^2nm}=-\frac4{\pi^2n}.
\]
This proves the asserted values of $C_B$ and $\beta_B$, because
$\int_0^1G_B=I_B+I_2/\pi$ and $a=2/\pi$.
The remaining error is bounded by
\[
 \frac{12}{m^2}+\frac{25}{3nm^2}+\frac{17}{4nm}<\frac{17}{n^2}.
\]
For the last inequality put $y=n/m<1$.  After multiplication by
$n^2$ the left side is
$(11/3)y^2+(151/12)y<65/4<17$.

\paragraph{The smooth first correction has a strict sign.}
Termwise differentiation of \eqref{eq:asym-details-B-series} gives
$B'(t)<0$.  Thus $0<B(t)+1/\pi<2/\pi$ for $0<t<1$.
Isolating the first pole in
\eqref{eq:asym-details-cot-fraction} gives
\[
 B(t)=\frac1{\pi(1+t)}-\frac{2t}{\pi}
                         \sum_{k\geq2}\frac1{k^2-t^2}.
\]
The last sum is greater than
$\sum_{k\geq2}k^{-2}>\int_2^\infty x^{-2}dx=1/2$, whence
\begin{equation}\label{eq:asym-details-B-sign-bound}
 B(t)+1/\pi<\frac{2-t^2}{\pi(1+t)}.
\end{equation}
The telescoping sum $\sum_{k\geq1}(4k^2-1)^{-1}=1/2$ and the
same cotangent expansion show
\[
 \frac{W(t)}{1-t^2}=\frac2\pi+\frac{4t^2}{\pi}
    \sum_{k\geq1}\frac1{(4k^2-1)(4k^2-t^2)}.
\]
Bound $(4k^2-t^2)^{-1}$ by $(4-t^2)^{-1}$ to obtain
\[
 W(t)\leq\frac2\pi\frac{1-t^2}{1-t^2/4}.
\]
Together with \eqref{eq:asym-details-B-sign-bound}, this yields
\[
 0<G_B(t)<\frac2{\pi^2}
       \frac{(1-t)(2-t^2)}{1-t^2/4}
       <\frac4{\pi^2}(1-t),\qquad 0<t<1.
\]
Thus $0<\int_0^1G_B<2/\pi^2$, and
$\beta_B=-4/\pi^2+2\int_0^1G_B$ lies strictly between
$-a^2$ and $0$.  This completes the proof of
Proposition~\ref{prop:asym-two-expansions}.
